% ====================================================================
% ch3v22_bib.tex — Chapter 3(Si·혼합) 장별 참고문헌 (14종) — R1 자동 분할(cite 스캔 기반)
%   원장 = results/V1022_REFERENCE_LEDGER.md (V1 키만). 장 자기완결(리뷰 모음형) — 공통 키의 타 장 중복 수록은 의도(D22).
% ====================================================================
\begin{thebibliography}{99}
\bibitem{wen_huggins1981} C. J. Wen and R. A. Huggins, ``Chemical diffusion in intermediate phases in the lithium-silicon system,'' \emph{J. Solid State Chem.} \textbf{37}(3), 271--278 (1981). DOI: 10.1016/0022-4596(81)90487-4. [고온 평형 Li--Si 중간상 titration 원전.]
\bibitem{limthongkul2003} P. Limthongkul, Y.-I. Jang, N. J. Dudney, and Y.-M. Chiang, ``Electrochemically-driven solid-state amorphization in lithium-silicon alloys and implications for lithium storage,'' \emph{Acta Mater.} \textbf{51}(4), 1103 (2003). DOI: 10.1016/S1359-6454(02)00515-4.
\bibitem{li_dahn2007} J. Li and J. R. Dahn, ``An In Situ X-Ray Diffraction Study of the Reaction of Li with Crystalline Si,'' \emph{J. Electrochem. Soc.} \textbf{154}(3), A156--A161 (2007). DOI: 10.1149/1.2409862.
\bibitem{obrovac_christensen2004} M. N. Obrovac and L. Christensen, ``Structural Changes in Silicon Anodes during Lithium Insertion/Extraction,'' \emph{Electrochem. Solid-State Lett.} \textbf{7}(5), A93--A96 (2004). DOI: 10.1149/1.1652421. [Li$_{15}$Si$_4$ 준안정 결정화.]
\bibitem{chevrier_dahn2009} V. L. Chevrier and J. R. Dahn, ``First Principles Model of Amorphous Silicon Lithiation,'' \emph{J. Electrochem. Soc.} \textbf{156}(6), A454--A458 (2009). DOI: 10.1149/1.3111037. [비정질 경로의 경사 전위-조성 곡선.]
\bibitem{beaulieu2001} L. Y. Beaulieu, K. W. Eberman, R. L. Turner, L. J. Krause, and J. R. Dahn, ``Colossal Reversible Volume Changes in Lithium Alloys,'' \emph{Electrochem. Solid-State Lett.} \textbf{4}(9), A137 (2001). DOI: 10.1149/1.1388178.
\bibitem{sethuraman_stressevo2010} V. A. Sethuraman, M. J. Chon, M. Shimshak, V. Srinivasan, and P. R. Guduru, ``In situ measurements of stress evolution in silicon thin films during electrochemical lithiation and delithiation,'' \emph{J. Power Sources} \textbf{195}(15), 5062--5066 (2010). DOI: 10.1016/j.jpowsour.2010.02.013. [소성 유동 $\sim-1.75$ GPa.]
\bibitem{sethuraman_stresspot2010} V. A. Sethuraman, V. Srinivasan, A. F. Bower, and P. R. Guduru, ``In Situ Measurements of Stress-Potential Coupling in Lithiated Silicon,'' \emph{J. Electrochem. Soc.} \textbf{157}(11) (2010). DOI: 10.1149/1.3489378. [응력--전위 결합 $\sim100$--$120$ mV/GPa --- 쪽 범위는 사용 시 Crossref 최종 대조.]
\bibitem{liu_sizefracture2012} X. H. Liu, L. Zhong, S. Huang, S. X. Mao, T. Zhu, and J. Y. Huang, ``Size-Dependent Fracture of Silicon Nanoparticles During Lithiation,'' \emph{ACS Nano} \textbf{6}(2), 1522--1531 (2012). DOI: 10.1021/nn204476h. [임계 $\sim150$ nm.]
\bibitem{obrovac_chevrier2014} M. N. Obrovac and V. L. Chevrier, ``Alloy Negative Electrodes for Li-Ion Batteries,'' \emph{Chem. Rev.} \textbf{114}(23), 11444--11502 (2014). DOI: 10.1021/cr500207g. [합금 음극 총람 리뷰, tier B.]
\bibitem{verbrugge_lisi2016} M. W. Verbrugge, D. R. Baker, and X. Xiao, ``Formulation for the Treatment of Multiple Electrochemical Reactions and Associated Speciation for the Lithium-Silicon Electrode,'' \emph{J. Electrochem. Soc.} \textbf{163}(2), A262--A271 (2016). DOI: 10.1149/2.0581602jes. [다반응$\cdot$speciation --- 전하 보존 구조의 Si 실증(MSMR 계보).]
\bibitem{jiang_sihys2020} Y. Jiang, G. Offer, J. Jiang, M. Marinescu, and H. Wang, ``Voltage Hysteresis Model for Silicon Electrodes for Lithium Ion Batteries, Including Multi-Step Phase Transformations, Crystallization and Amorphization,'' \emph{J. Electrochem. Soc.} \textbf{167}(13), 130533 (2020). DOI: 10.1149/1945-7111/abbbba.
\bibitem{larchecahn1973} F. Larch\'e and J. W. Cahn, ``A linear theory of thermochemical equilibrium of solids under stress,'' \emph{Acta Metall.} \textbf{21}(8), 1051--1063 (1973). DOI: 10.1016/0001-6160(73)90021-7. [응력--조성 결합 화학퍼텐셜의 이론 원전.]
\bibitem{koebbing2024} L. K\"obbing, A. Latz, and B. Horstmann, ``Voltage Hysteresis of Silicon Nanoparticles: Chemo-Mechanical Particle-SEI Model,'' \emph{Adv. Funct. Mater.}, 2308818. DOI: 10.1002/adfm.202308818. [SEI 점탄소성 소산 히스 모델 --- 인쇄 권$\cdot$호는 사용 시 Crossref 최종 대조.]
\end{thebibliography}
